\documentclass[11pt,a4paper]{article}
% preamble.tex — shared setup for human-verification documents
% Compile each document with:  tectonic <File>.tex

\usepackage[T1]{fontenc}
\usepackage[utf8]{inputenc}
\usepackage{lmodern}
\usepackage{amssymb,amsmath,amsthm}
\usepackage{mathtools}
\usepackage[dvipsnames]{xcolor}
\usepackage{listings}
\usepackage{hyperref}
\usepackage[margin=2.5cm]{geometry}
\usepackage{parskip}

% ── Colours ──────────────────────────────────────────────────────────
\definecolor{leanBg}{HTML}{F7F7F0}
\definecolor{haskellBg}{HTML}{F0F4F7}
\definecolor{keyword}{HTML}{7B3294}
\definecolor{comment}{HTML}{5A8C32}
\definecolor{noteFg}{HTML}{B35900}

% ── Lean 4 listings language ─────────────────────────────────────────
% NOTE: Lean listing blocks use ASCII approximations of unicode symbols
% (e.g. -> for arrow, o. for compose).  The literate key substitutes
% these with proper math-mode renderings in the PDF.
\lstdefinelanguage{Lean4}{
  morekeywords={def,theorem,structure,where,namespace,end,import,fun,
                let,have,show,by,exact,intro,cases,rfl,subst,obtain,
                funext,simp,rw,apply,variable,Type,Prop,instance,
                section,open,class,extends,deriving,sorry,noncomputable},
  sensitive=true,
  morecomment=[l]{--},
  morecomment=[n]{/-}{-/},
  morestring=[b]",
}

% ── Lean listing style ───────────────────────────────────────────────
\lstdefinestyle{lean}{
  language=Lean4,
  backgroundcolor=\color{leanBg},
  basicstyle=\small\ttfamily,
  keywordstyle=\color{keyword}\bfseries,
  commentstyle=\color{comment}\itshape,
  breaklines=true,
  frame=single,
  framerule=0.4pt,
  xleftmargin=4pt,
  xrightmargin=4pt,
  aboveskip=6pt,
  belowskip=4pt,
  literate=
    {->}{{$\to$}}2
    {=>}{{$\Rightarrow$}}2
    {<|}{{$\langle$}}1
    {|>}{{$\rangle$}}1
    {==>}{{$\Longrightarrow$}}3
    { o. }{{$\;\circ\;$}}3
    {forall}{{$\forall$}}1
    {:=}{{$\,:=\,$}}2
    {!=}{{$\neq$}}2
    {<=}{{$\leq$}}2
    {>=}{{$\geq$}}2,
}

% ── Haskell listing style ───────────────────────────────────────────
\lstdefinestyle{haskell}{
  language=Haskell,
  backgroundcolor=\color{haskellBg},
  basicstyle=\small\ttfamily,
  keywordstyle=\color{keyword}\bfseries,
  commentstyle=\color{comment}\itshape,
  breaklines=true,
  frame=single,
  framerule=0.4pt,
  xleftmargin=4pt,
  xrightmargin=4pt,
  aboveskip=6pt,
  belowskip=4pt,
  literate=
    {->}{{$\to$}}2
    {=>}{{$\Rightarrow$}}2,
}

% ── Theorem environments ────────────────────────────────────────────
\theoremstyle{definition}
\newtheorem{defn}{Definition}[section]
\newtheorem{thm}[defn]{Theorem}
\newtheorem{prop}[defn]{Proposition}

\theoremstyle{remark}
\newtheorem*{fidelitynote}{Fidelity note}

% ── Convenience macros ───────────────────────────────────────────────
\newcommand{\lean}[1]{\texttt{#1}}      % inline Lean identifier
\newcommand{\hask}[1]{\texttt{#1}}      % inline Haskell identifier
\newcommand{\Ext}{\mathrm{Ext}}
\newcommand{\Pos}{\mathrm{Pos}}
\newcommand{\Shape}{\mathrm{Shape}}


\title{Human Verification: \texttt{Containers.CoKleisli}\\[4pt]
  \large CoKleisli Category and Canonical Strength}
\author{Machine-generated from \texttt{lean/CoKleisli.lean} for review by Neil Ghani}
\date{July 2026}

\begin{document}
\maketitle

\begin{abstract}
This document presents every definition and theorem \emph{statement} from the
Lean~4 module \lean{Containers.CoKleisli} in three parallel forms: the original
Lean, a readable Haskell rendering, and precise mathematical English.  Proofs
are omitted --- they are machine-checked by Lean.  The purpose is to let a
human verifier confirm that the formalisation captures the intended mathematics.
\end{abstract}

\tableofcontents
\bigskip

%===================================================================
\section{Overview}
%===================================================================

This file develops the coKleisli structure and canonical strength for directed
container comonads, extending \lean{Containers.Directed}.

A \textbf{directed container} $D = (S, P, o, \downarrow, \oplus)$ gives rise to a
comonad on its extension $\mathbf{Ext}\,D$ (a comonad on $\mathbf{Type}$).  The
coKleisli category of this comonad has as morphisms $X \to Y$ the functions
$\mathbf{Ext}\,D\,X \to Y$, composed via the \emph{extension} operation.  This file
defines \lean{extend}, \lean{cokleisli\_comp}, and the canonical strength
$\sigma : W(A \times B) \to A \times W(B)$, and proves the categorical laws.

Everything is \textbf{Type}-level, Lean~4 core, no Mathlib.

%===================================================================
\section{CoKleisli extension}
%===================================================================

% --- 1. extend ----------------------------------------------------
\subsection{Definition: \lean{extend}}

\begin{defn}[CoKleisli extension]
Given $f : \mathbf{Ext}\,D\,X \to Y$, the \textbf{coKleisli extension}
$\mathtt{extend}\,f : \mathbf{Ext}\,D\,X \to \mathbf{Ext}\,D\,Y$
applies $f$ at every position via comultiplication:
\[
  \mathtt{extend}\,f\,\langle s, v \rangle
  \;=\;
  \bigl\langle s,\;
    \lambda p.\; f\bigl\langle s \mathop{\downarrow} p,\;
                    \lambda q.\; v(p \oplus q)\bigr\rangle
  \bigr\rangle.
\]
Equivalently, $\mathtt{extend}\,f = \mathbf{Ext}\,f \circ \delta$, where
$\delta$ is the comonad comultiplication.
\end{defn}

\paragraph{Lean.}
\begin{lstlisting}[style=lean]
def extend {X Y : Type} (f : Ext D.toContainer X -> Y) :
    Ext D.toContainer X -> Ext D.toContainer Y :=
  fun p => <|p.1, fun q => f <|D.sub p.1 q, fun r => p.2 (D.shift p.1 q r)|>|>
\end{lstlisting}

\paragraph{Haskell.}
\begin{lstlisting}[style=haskell]
-- In the coKleisli category of a comonad W, extend lifts
-- f :: W x -> y  to  extend f :: W x -> W y.
-- For a directed container this is:
extend :: DirectedContainer s p -> (Ext s p x -> y) -> Ext s p x -> Ext s p y
extend d f (Ext s v) = Ext s (\q -> f (Ext (sub d s q) (\r -> v (shift d s q r))))
\end{lstlisting}


% --- 2. extend_eq_map_comult --------------------------------------
\subsection{Theorem: \lean{extend\_eq\_map\_comult}}

\begin{prop}[Extension as map after comult]
For any $f : \mathbf{Ext}\,D\,X \to Y$,
\[
  \mathtt{extend}\,f \;=\; \mathbf{Ext}\mathord{.}\mathtt{map}\,f \circ \delta,
\]
where $\delta = \mathtt{comult}$ is the comonad comultiplication.  This holds
definitionally (by \lean{rfl}).
\end{prop}

\paragraph{Lean.}
\begin{lstlisting}[style=lean]
theorem extend_eq_map_comult {X Y : Type} (f : Ext D.toContainer X -> Y) :
    D.extend f = Ext.map f o. D.comult
\end{lstlisting}

\paragraph{Haskell.}
\begin{lstlisting}[style=haskell]
-- extend f = fmap f . comult
-- This is the standard relationship between extend and (>>=) in a comonad.
\end{lstlisting}


% --- 3. extend_counit ---------------------------------------------
\subsection{Theorem: \lean{extend\_counit}}

\begin{thm}[\lean{extend} of the counit is the identity]
\[
  \mathtt{extend}\,\varepsilon \;=\; \mathrm{id}.
\]
This is the \emph{right unit law} of the coKleisli category.
\end{thm}

\paragraph{Lean.}
\begin{lstlisting}[style=lean]
theorem extend_counit {X : Type} :
    D.extend (D.counit (X := X)) = id
\end{lstlisting}

\paragraph{Haskell.}
\begin{lstlisting}[style=haskell]
-- Comonad law (right unit):
-- extend extract = id
\end{lstlisting}


% --- 4. counit_extend ---------------------------------------------
\subsection{Theorem: \lean{counit\_extend}}

\begin{thm}[Counit is a left inverse of extension]
For every $f : \mathbf{Ext}\,D\,X \to Y$,
\[
  \varepsilon \circ \mathtt{extend}\,f \;=\; f.
\]
This characterises $\mathtt{extend}\,f$ as the unique map $h$ satisfying
$\varepsilon \circ h = f$ (the universal property of extension).
\end{thm}

\paragraph{Lean.}
\begin{lstlisting}[style=lean]
theorem counit_extend {X Y : Type} (f : Ext D.toContainer X -> Y) :
    D.counit o. D.extend f = f
\end{lstlisting}

\paragraph{Haskell.}
\begin{lstlisting}[style=haskell]
-- Comonad law (left unit):
-- extract . extend f = f
\end{lstlisting}


% --- 5. extend_extend ---------------------------------------------
\subsection{Theorem: \lean{extend\_extend}}

\begin{thm}[Extension respects coKleisli composition]
For $f : \mathbf{Ext}\,D\,X \to Y$ and $g : \mathbf{Ext}\,D\,Y \to Z$,
\[
  \mathtt{extend}(g \circ \mathtt{extend}\,f)
  \;=\;
  \mathtt{extend}\,g \;\circ\; \mathtt{extend}\,f.
\]
This is the \emph{associativity law} of the coKleisli category; it follows
from coassociativity of the comonad.
\end{thm}

\paragraph{Lean.}
\begin{lstlisting}[style=lean]
theorem extend_extend {X Y Z : Type}
    (f : Ext D.toContainer X -> Y)
    (g : Ext D.toContainer Y -> Z) :
    D.extend (g o. D.extend f) = D.extend g o. D.extend f
\end{lstlisting}

\paragraph{Haskell.}
\begin{lstlisting}[style=haskell]
-- Comonad law (associativity / fusion):
-- extend (g . extend f) = extend g . extend f
\end{lstlisting}


%===================================================================
\section{CoKleisli composition}
%===================================================================

% --- 6. cokleisli_comp --------------------------------------------
\subsection{Definition: \lean{cokleisli\_comp}}

\begin{defn}[CoKleisli composition]
Given coKleisli morphisms $f : \mathbf{Ext}\,D\,X \to Y$ and
$g : \mathbf{Ext}\,D\,Y \to Z$, their \textbf{coKleisli composite} is
\[
  g \mathbin{.*} f \;:=\; g \circ \mathtt{extend}\,f
  \;:\; \mathbf{Ext}\,D\,X \to Z.
\]
\end{defn}

\paragraph{Lean.}
\begin{lstlisting}[style=lean]
def cokleisli_comp {X Y Z : Type}
    (f : Ext D.toContainer X -> Y)
    (g : Ext D.toContainer Y -> Z) :
    Ext D.toContainer X -> Z :=
  g o. D.extend f
\end{lstlisting}

\paragraph{Haskell.}
\begin{lstlisting}[style=haskell]
-- CoKleisli composition  (=>>) in Control.Comonad:
-- (=>>) :: (w a -> b) -> (w b -> c) -> w a -> c
cokleisliComp :: (Ext s p x -> y) -> (Ext s p y -> z) -> Ext s p x -> z
cokleisliComp f g = g . extend f
\end{lstlisting}


% --- 7. cokleisli_comp_counit -------------------------------------
\subsection{Theorem: \lean{cokleisli\_comp\_counit}}

\begin{thm}[Right unit of CoKleisli composition]
For every $f : \mathbf{Ext}\,D\,X \to Y$,
\[
  f \mathbin{.*} \varepsilon \;=\; f.
\]
The counit $\varepsilon$ is a right identity for coKleisli composition.
\end{thm}

\paragraph{Lean.}
\begin{lstlisting}[style=lean]
theorem cokleisli_comp_counit {X Y : Type}
    (f : Ext D.toContainer X -> Y) :
    D.cokleisli_comp (D.counit) f = f
\end{lstlisting}

\paragraph{Haskell.}
\begin{lstlisting}[style=haskell]
-- Right unit:  f =>> extract = f
\end{lstlisting}


% --- 8. cokleisli_counit_comp -------------------------------------
\subsection{Theorem: \lean{cokleisli\_counit\_comp}}

\begin{thm}[Left unit of CoKleisli composition]
For every $f : \mathbf{Ext}\,D\,X \to Y$,
\[
  \varepsilon \mathbin{.*} f \;=\; f.
\]
The counit $\varepsilon$ is a left identity for coKleisli composition.
\end{thm}

\paragraph{Lean.}
\begin{lstlisting}[style=lean]
theorem cokleisli_counit_comp {X Y : Type}
    (f : Ext D.toContainer X -> Y) :
    D.cokleisli_comp f D.counit = f
\end{lstlisting}

\paragraph{Haskell.}
\begin{lstlisting}[style=haskell]
-- Left unit:  extract =>> f = f
\end{lstlisting}


% --- 9. cokleisli_assoc -------------------------------------------
\subsection{Theorem: \lean{cokleisli\_assoc}}

\begin{thm}[Associativity of CoKleisli composition]
For $f : \mathbf{Ext}\,D\,X \to Y$, $g : \mathbf{Ext}\,D\,Y \to Z$,
$h : \mathbf{Ext}\,D\,Z \to W$,
\[
  (h \mathbin{.*} g) \mathbin{.*} f \;=\; h \mathbin{.*} (g \mathbin{.*} f).
\]
\end{thm}

\paragraph{Lean.}
\begin{lstlisting}[style=lean]
theorem cokleisli_assoc {X Y Z W : Type}
    (f : Ext D.toContainer X -> Y)
    (g : Ext D.toContainer Y -> Z)
    (h : Ext D.toContainer Z -> W) :
    D.cokleisli_comp (D.cokleisli_comp f g) h
      = D.cokleisli_comp f (D.cokleisli_comp g h)
\end{lstlisting}

\paragraph{Haskell.}
\begin{lstlisting}[style=haskell]
-- Associativity:  (f =>> g) =>> h = f =>> (g =>> h)
\end{lstlisting}


%===================================================================
\section{Canonical strength}
%===================================================================

% --- 10. strength -------------------------------------------------
\subsection{Definition: \lean{strength}}

\begin{defn}[Canonical strength]
The \textbf{canonical strength} of the directed container comonad is a natural
transformation
\[
  \sigma_{A,B} : \mathbf{Ext}\,D\,(A \times B) \;\longrightarrow\;
                 A \times \mathbf{Ext}\,D\,B,
\]
defined by
\[
  \sigma_{A,B}\langle s, v \rangle
  \;=\;
  \bigl(\pi_1(v(o(s))),\;
        \langle s, \pi_2 \circ v \rangle\bigr).
\]
The first component extracts the $A$-value at the root; the second keeps the
shape and strips each labelled pair to its $B$-component.
\end{defn}

\paragraph{Lean.}
\begin{lstlisting}[style=lean]
def strength (A B : Type) :
    Ext D.toContainer (A × B) -> A × Ext D.toContainer B :=
  fun p => <|(p.2 (D.root p.1)).1, <|p.1, fun q => (p.2 q).2|>|>
\end{lstlisting}

\paragraph{Haskell.}
\begin{lstlisting}[style=haskell]
-- Canonical comonad strength:  sigma :: w (a, b) -> (a, w b)
-- For a directed container:
strength :: DirectedContainer s p -> Ext s p (a, b) -> (a, Ext s p b)
strength d (Ext s v) = (fst (v (root d s)), Ext s (snd . v))
\end{lstlisting}

\begin{fidelitynote}
Every comonad on a cartesian category is canonically strong.  Here the
$A$-component is extracted at the \emph{root} position $o(s)$, which is
well-typed because the counit extracts at the root.  The Lean source has a
minor syntactic issue in line~166 (a missing opening angle bracket); the
intent is clear from context and from the surrounding docstring.
\end{fidelitynote}


% --- 11. strength_natural_left ------------------------------------
\subsection{Theorem: \lean{strength\_natural\_left}}

\begin{thm}[Strength naturality in $A$]
For every $f : A \to A'$,
\[
  (f \times \mathrm{id}) \circ \sigma_{A,B}
  \;=\;
  \sigma_{A',B} \circ \mathbf{Ext}\,D\,(f \times \mathrm{id}).
\]
\end{thm}

\paragraph{Lean.}
\begin{lstlisting}[style=lean]
theorem strength_natural_left {A A' B : Type} (f : A -> A') :
    (fun p : A × Ext D.toContainer B => (f p.1, p.2)) o. D.strength A B
      = D.strength A' B o. Ext.map (fun ab : A × B => (f ab.1, ab.2))
\end{lstlisting}

\paragraph{Haskell.}
\begin{lstlisting}[style=haskell]
-- Left naturality:
-- first f . strength = strength . fmap (first f)
-- where first f (a, b) = (f a, b)
\end{lstlisting}


% --- 12. strength_natural_right -----------------------------------
\subsection{Theorem: \lean{strength\_natural\_right}}

\begin{thm}[Strength naturality in $B$]
For every $g : B \to B'$,
\[
  (\mathrm{id} \times \mathbf{Ext}\,D\,g) \circ \sigma_{A,B}
  \;=\;
  \sigma_{A,B'} \circ \mathbf{Ext}\,D\,(\mathrm{id} \times g).
\]
\end{thm}

\paragraph{Lean.}
\begin{lstlisting}[style=lean]
theorem strength_natural_right {A B B' : Type} (g : B -> B') :
    (fun p : A × Ext D.toContainer B => (p.1, Ext.map g p.2)) o. D.strength A B
      = D.strength A B' o. Ext.map (fun ab : A × B => (ab.1, g ab.2))
\end{lstlisting}

\paragraph{Haskell.}
\begin{lstlisting}[style=haskell]
-- Right naturality:
-- second (fmap g) . strength = strength . fmap (second g)
-- where second g (a, b) = (a, g b)
\end{lstlisting}


% --- 13. strength_counit ------------------------------------------
\subsection{Theorem: \lean{strength\_counit}}

\begin{thm}[Strength--counit coherence]
\[
  (\mathrm{id}_A \times \varepsilon_B) \circ \sigma_{A,B}
  \;=\;
  \varepsilon_{A \times B}.
\]
Extracting from the $B$-component of the strength gives the same result as
extracting directly from the $A \times B$-tree.
\end{thm}

\paragraph{Lean.}
\begin{lstlisting}[style=lean]
theorem strength_counit {A B : Type} :
    (fun p : A × Ext D.toContainer B => (p.1, D.counit p.2)) o. D.strength A B
      = D.counit (X := A × B)
\end{lstlisting}

\paragraph{Haskell.}
\begin{lstlisting}[style=haskell]
-- Strength-counit coherence:
-- second extract . strength = extract
-- i.e.  (id *** extract) . sigma = extract
\end{lstlisting}


%===================================================================
\section{Translation conventions}
%===================================================================

The following conventions apply throughout this document.

\begin{enumerate}
\item \textbf{Directed container notation.}  A directed container
  $D = (S, P, o, \downarrow, \oplus)$ has shapes $S$, positions $P$,
  root $o : S \to P$, sub-shape $(\downarrow) : S \to P(s) \to S$, and
  shift $(\oplus) : (s : S) \to P(s) \to P(s \downarrow p) \to P(s)$.
  The laws D1--D5 from \lean{Containers.Directed} are used implicitly in the
  proofs.

\item \textbf{CoKleisli = Kleisli of the opposite monad.}  The coKleisli
  category of a comonad $W$ has objects the types and morphisms $X \to Y$ the
  functions $W(X) \to Y$; identity is $\varepsilon$ and composition is
  $\mathtt{cokleisli\_comp}$.

\item \textbf{Haskell comonad vocabulary.}  Haskell's \hask{Control.Comonad}
  uses \hask{extract} for $\varepsilon$ and \hask{extend} for the coKleisli
  lift; the operator \hask{=>>} is coKleisli composition in the left-to-right
  order.

\item \textbf{Proof bodies omitted.}  Only definitions and theorem
  \emph{statements} are translated.  Lean checks the proofs; the human
  verifier checks that the \emph{statements} match the intended mathematics.
\end{enumerate}

\end{document}
